\documentclass[10pt,conference,compsocconf]{IEEEtran}

\usepackage[T2A]{fontenc}
\usepackage[utf8]{inputenc}
\usepackage[english,russian]{babel}

\usepackage[dvipsnames]{xcolor}


\usepackage{cite}
\makeatletter
\renewcommand{\@IEEEsectpunct}{.\ }
\makeatother
\usepackage{graphicx}    % For figure environment
\usepackage{booktabs}
\usepackage{multirow}
\usepackage{tabularx}
\usepackage{amsmath}
\usepackage{amsfonts}
\usepackage{pdfrender}
\usepackage[hidelinks,unicode]{hyperref}
\newcommand{\bestmetric}[1]{\textpdfrender{TextRenderingMode=FillStroke,LineWidth=0.25pt}{\usefont{OT1}{phv}{b}{n}#1}}


\begin{document}
\title{Аудиовизуальное разделение речи двух дикторов}

\author{
  Андрей Теплов\\
  \textit{Faculty of Computer Science, HSE, Moscow}\\
  \texttt{arteplov@edu.hse.ru}
  \and
  Никита Романенко\\
  \textit{Faculty of Computer Science, HSE, Moscow}\\
  \texttt{ndromanenko@hse.edu.ru}
}

\date{\today}

\maketitle

\begin{abstract}
В работе рассматривается задача разделения речи двух дикторов по записи их разговора. Мы исследуем, как на качество восстановления отдельных голосов влияют выбор архитектуры и использование видеозаписей говорящих. Для разделения речи только по звуку мы сравнили Conv-TasNet, Dual-Path RNN (DPRNN) и SepFormer. Для Conv-TasNet и DPRNN мы также проверили свёрточный видеоэнкодер Conv3D и покадровый Vision Transformer (ViT). На валидационной выборке DPRNN превзошла Conv-TasNet как при работе только со звуком, так и с аудиовизуальными экспериментами. Лучший результат среди рассмотренных моделей показала DPRNN с ViT.
\end{abstract}

\section{Введение}\label{sec:introduction}

Когда два человека говорят одновременно, их голоса накладываются в одноканальной записи. Задача разделения речи состоит в том, чтобы восстановить отдельную дорожку каждого диктора даже на участках перекрытия. Мы сравниваем Conv-TasNet, DPRNN и SepFormer при работе только со звуком. Для Conv-TasNet и DPRNN также проверяем, помогают ли синхронные видеозаписи движений губ разделить голоса и как результат зависит от видеоэнкодера и способа объединения признаков.


\section{Обзор существующих подходов}\label{sec:related_work}

В Conv-TasNet~\cite{luo2019convtasnet} энкодер преобразует звуковую волну в последовательность признаков, а сеть разделения оценивает маску каждого источника. Последовательные свёрточные блоки используют растущие коэффициенты расширения, поэтому при том же размере ядра каждый следующий слой охватывает более длинный фрагмент записи. Так при разделении наложившихся голосов модель учитывает и соседние звуки.

В DPRNN~\cite{luo2020dprnn} временной контекст моделируется на 2 уровнях. Последовательность делится на перекрывающиеся сегменты. Рекуррентный слой внутри сегмента обрабатывает локальные изменения речи, а слой между сегментами связывает одинаковые позиции в разных частях записи. SepFormer~\cite{subakan2021sepformer} сохраняет это разбиение, но использует блоки attention~\cite{vaswani2017attention}. Attention определяет, какие позиции внутри сегмента и какие другие сегменты наиболее важны для текущего предсказания.

В аудиовизуальном разделении дополнительным признаком служит связь между артикуляцией и звуком. Авторы статьи~\cite{ephrat2018looking} показали, что по последовательности изображений лица можно выделить из смеси голос выбранного человека. В нашей работе модели доступны записи области рта обоих дикторов. Это позволяет проверить, помогают ли визуальные признаки одновременно восстановить обе дорожки.

\section{Методология}\label{sec:methodology}

\subsection{Разделение речи}
Задача заключалась в восстановлении отдельных речевых сигналов из одноканальной смеси. Для двух дикторов смесь представлялась как $x(t)=s_1(t)+s_2(t)$, где $s_1(t)$ и $s_2(t)$ --- исходные сигналы дикторов. По смеси $x(t)$ модель оценивала дорожки $\hat{s}_1(t)$ и $\hat{s}_2(t)$.

Разделение выполнялось во временной области. Аудиоэнкодер преобразовывал смесь в последовательность признаков $Z=E(x)\in\mathbb{R}^{N\times T}$, где $N$ --- число каналов представления, а $T$ - число временных шагов после кодирования.

Для каждого диктора $i\in\{1,2\}$ сеть разделения оценивала маску $M_i\in[0,1]^{N\times T}$, выделяющую его вклад в представление смеси. Маску поэлементно умножали на $Z$, после чего общий для обоих дикторов декодер восстанавливал сигнал $\hat{s}_i=D(M_i\odot Z)$. В аудиовизуальных вариантах при оценке масок дополнительно использовались признаки видеозаписей обоих дикторов.

\subsection{Визуальные признаки и их объединение со звуком}
В варианте с Conv3D первая свёртка охватывала 5 соседних кадров, затем пространственные признаки обрабатывались покадрово и объединялись временной свёрткой. В варианте ViT~\cite{dosovitskiy2021vit} attention связывал участки внутри каждого кадра. Усреднение по участкам и линейная проекция давали 32 признака на кадр, как и в Conv3D.

Чтобы учитывать движения губ при оценке масок, модель объединяла признаки видеоэнкодера с признаками звуковой смеси. В Add 64 видеоканала обоих дикторов проецировали в 128 каналов, линейно интерполировали до длины $T$ и прибавляли к аудиопризнакам. В Concat-MLP после интерполяции соединяли 64 видеоканала со 128 аудиоканалами. Их сочетание на каждом шаге определяла обучаемая сеть $192\!\rightarrow\!128\!\rightarrow\!128$ с ReLU между слоями. В Conv-TasNet признаки объединяли после 1-й группы свёрточных блоков, в DPRNN - перед разбиением на сегменты.

\subsection{Режимы обучения}

\textit{Только аудио.} Conv-TasNet, DPRNN и SepFormer восстанавливали 2 речевые дорожки по одноканальной смеси. Функцию потерь PIT SI-SNR~\cite{kolbaek2017upit} вычисляли для каждой смеси по лучшему соответствию между выходами модели и эталонами
\[
\mathcal{L}_{\mathrm{sep}}=-\max_{\pi\in S_2}\frac{1}{2}\sum_{i=1}^{2}\text{SI-SNR}(\hat{s}_i,s_{\pi(i)}),
\]
где $\hat{s}_i$ --- восстановленная дорожка, $s_{\pi(i)}$ --- соответствующая ей чистая дорожка, а $S_2$ содержит 2 перестановки дикторов. При обучении потери усредняли по смесям в батче.

\textit{Аудио и видео.} Conv-TasNet и DPRNN дополнительно получали синхронные видео области рта обоих дикторов. Модель использовала визуальные признаки при оценке масок, а параметры обучали по той же функции потерь $\mathcal{L}_{\mathrm{sep}}$.

\textit{Аудио и видео с двумя функциями потерь.} Для DPRNN + ViT к PIT SI-SNR добавили задачу сопоставления видео с голосом. Учителем для ViT служил только энкодер из чекпоинта Conv-TasNet. Его веса заморозили, чтобы признаки чистой речи каждого диктора оставались фиксированным эталоном. Для каждого видео сравнивали проекцию признаков ViT с обоими эталонами по среднему покадровому косинусному сходству. Cross энтропия $\mathcal{L}_{\mathrm{CE}}$ с температурой $\tau=0.1$ поощряла большее сходство с дорожкой диктора в кадре. Итоговая функция потерь имела вид $\mathcal{L}_{\mathrm{sep}}+\lambda\mathcal{L}_{\mathrm{CE}}$. Учитель требовался только для вычисления CE. При инференсе модель разделяла смесь по звуку и видео самостоятельно.

\section{Постановка экспериментов}\label{sec:experimental_setup}

В экспериментах использовали набор из 28\,000 одноканальных смесей речи двух дикторов. Каждая смесь длится 2~с (32\,000 отсчётов при 16~кГц) и сопровождается видео области рта каждого диктора. Обучающая выборка содержит 20\,000 смесей, валидационная - 5\,000, тестовая - 3\,000. Чистые речевые дорожки доступны для обучения и валидации.

\subsection{Условия сравнения}
В таблицах приведены метрики на валидационной выборке. SI-SNR характеризует точность восстановления сигнала, SI-SNRi и SDRi~--- улучшение относительно смеси, PESQ --- воспринимаемое качество речи, STOI --- её разборчивость. Для каждой метрики жирным выделен максимум в пределах таблицы, определённый до округления.

\subsection{Конфигурации видеоэнкодеров}

На вход каждого энкодера подавали 50 полутоновых кадров области рта размером $96 \times 96$, на выходе получали 32 признака на кадр. В табл.~\ref{tab:visual_encoders} указано число параметров одного видеоэнкодера с общими для 2 дикторов весами.

\makeatletter
\newenvironment{experimenttable}{%
  \par\addvspace{\medskipamount}\noindent\begin{minipage}{\linewidth}%
  \def\@captype{table}\centering\footnotesize\setlength{\tabcolsep}{3pt}%
}{\end{minipage}\par\nobreak\smallskip}
\makeatother

\begin{experimenttable}
    \caption{Гиперпараметры видеоэнкодеров}
    \label{tab:visual_encoders}
    \begin{tabularx}{\linewidth}{@{}lXr@{}}
        \toprule
        Энкодер & Параметр & Значение \\
        \midrule
        Conv3D & \texttt{Conv3d.kernel\_size} & $(5,5,5)$ \\
               & \texttt{Conv3d.stride} & $(1,4,4)$ \\
               & \texttt{Conv3d.out\_channels} & 16 \\
               & \texttt{Conv2d.out\_channels} & 32 \\
               & \texttt{Conv1d.kernel\_size} & 5 \\
               & \textbf{Всего параметров} & \textbf{21\,056} \\
        \midrule
        ViT    & \texttt{patch\_size} & 8 \\
               & \texttt{num\_patches} & 144 \\
               & \texttt{hidden\_size} & 128 \\
               & \texttt{num\_layers} & 4 \\
               & \texttt{num\_heads} & 4 \\
               & \texttt{dim\_feedforward} & 512 \\
               & \texttt{dropout} & 0.1 \\
               & \textbf{Всего параметров} & \textbf{824\,224} \\
        \bottomrule
    \end{tabularx}
\end{experimenttable}

\section{Результаты экспериментов}\label{sec:results}

\subsection{Эксперименты с аудиомоделями}\label{sec:audio_experiments}

Все 3 модели начинаются с энкодера --- свёртки с $N$ фильтрами и окном $L$ отсчётов. Каждому окну длиной $L$ соответствует кадр с $N$ признаками. Модели обучали с нуля только по звуку с функцией потерь PIT SI-SNR. Для Conv-TasNet и DPRNN использовали Adam с LR $10^{-3}$, батч 32, AMP, ограничение нормы градиента 5, seed 42 и ReduceLROnPlateau по SI-SNRi на валидации. Эпоха это полный проход по 20\,000 семплов длительностью 2с. Метрики посчитаны на валидации из 5\,000 семплов. SI-SNR не приводится, так как совпадает с SI-SNRi до 0.01~дБ. SI-SNRi и SDRi указаны в дБ, параметры --- в миллионах, жирным выделен максимум таблицы.

\subsubsection{Conv-TasNet}

В базовой конфигурации энкодер содержал $N=256$ фильтров длиной $L=32$ отсчёта. Перед свёрточными блоками признаки сжимаются до $B=128$ каналов, внутри блока используется $H=256$ каналов и ядро $P=3$. Повтор состоит из $X=8$ блоков с растущем dilation, всего $R=4$ повтора (3.34\,M параметров). Запуски обучали до 100 эпох с ранней остановкой (табл.~\ref{tab:ctn_ablation}).

\begin{experimenttable}
    \caption{Абляции Conv-TasNet, 100 эпох}\label{tab:ctn_ablation}
    \setlength{\tabcolsep}{1.8pt}
    \begin{tabularx}{\linewidth}{@{}>{\raggedright\arraybackslash}Xrrrrr@{}}
        \toprule
Изменение & Парам. & SI-SNRi & SDRi & PESQ & STOI \\
        \midrule
Базовая  & 3.34 & 10.39 & 10.80 & 1.82 & 0.876 \\
$L=16$   & 3.34 & 10.70 & 11.11 & \bestmetric{1.93} & 0.880 \\
$H=512$  & 6.56 & \bestmetric{10.75} & \bestmetric{11.16} & 1.92 & \bestmetric{0.881} \\
$N=512$  & 3.46 & 10.13 & 10.54 & 1.86 & 0.872 \\
$X=6$    & 2.54 & 10.17 & 10.59 & 1.87 & 0.873 \\
$R=2$    & 1.73 & 9.09  & 9.49  & 1.75 & 0.856 \\
        \bottomrule
    \end{tabularx}
\end{experimenttable}

Удвоение числа каналов в свёрточных блоках до $H=512$ повысило SI-SNRi на $0.36$, а уменьшение длины ядра вдвое до $L=16$ дало прирост $0.31$. В первом случае большая размерность, вероятно, позволяет лучше описывать паттерны в речи дикторов. Во втором энкодер делит звук на более короткие фрагменты. Речь быстро меняется во времени, и короткие фрагменты, вероятно, позволяют лучше уловить эти изменения и точнее определить, где чей голос. Меньшее число блоков ухудшило результат (на $1.30$ при $R=2$ и на $0.21$ при $X=6$), так как модель видит более короткий отрезок звука. Увеличение $N$ не дало прироста. Одно из возможных объяснений состоит в том, что перед сетью разделения признаки сжимаются до 128 каналов.

\subsubsection{DPRNN}

В базовой конфигурации использовали $N=128$ фильтров длиной $L=16$ отсчётов. Последовательность кадров делится на чанки длиной $K=100$, и 6 DPRNN-блоков поочерёдно применяют LSTM с размерностью $H=128$ внутри чанков и между ними (3.61\,M параметров). Все запуски в табл.~\ref{tab:dprnn_ablation} сравниваются на бюджете 30 эпох.

\begin{experimenttable}
    \caption{Абляции DPRNN, 30 эпох}\label{tab:dprnn_ablation}
    \setlength{\tabcolsep}{1.8pt}
    \begin{tabularx}{\linewidth}{@{}>{\raggedright\arraybackslash}Xrrrrr@{}}
        \toprule
Изменение & Парам. & SI-SNRi & SDRi & PESQ & STOI \\
        \midrule
Базовая          & 3.61  & 9.83 & 10.28 & 1.78 & 0.861 \\
$L=8$, батч 16   & 3.60  & \bestmetric{10.17} & \bestmetric{10.61} & \bestmetric{1.84} & \bestmetric{0.870} \\
$L=32$           & 3.61  & 8.15 & 8.57  & 1.61 & 0.834 \\
$K=50$           & 3.61  & 9.40 & 9.84  & 1.76 & 0.859 \\
$K=200$          & 3.61  & 9.45 & 9.89  & 1.75 & 0.856 \\
$H=64$           & 1.43  & 8.46 & 8.87  & 1.63 & 0.840 \\
$H=256$, батч 16 & 10.31 & 9.68 & 10.09 & 1.72 & 0.860 \\
8 блоков         & 4.79  & 9.61 & 10.06 & 1.78 & 0.857 \\
$N=64$           & 2.59  & 9.32 & 9.75  & 1.71 & 0.851 \\
$N=256$          & 5.68  & 9.03 & 9.47  & 1.71 & 0.850 \\
        \bottomrule
    \end{tabularx}
\end{experimenttable}

На DPRNN длина ядра энкодера влияла сильнее, чем на Conv-TasNet. При $L=8$ и батче 16 прирост SI-SNRi составил $0.34$, а при $L=32$ результат снизился на $1.68$. Лучшей оказалась длина чанка $K=100$, при которой обе RNN обрабатывают последовательности примерно одинаковой длины, как рекомендуют авторы DPRNN~\cite{luo2020dprnn}. Уменьшение $H$ до 64 снизило SI-SNRi на $1.37$, а уменьшение $N$ до 64 --- на $0.51$. Вероятно, при меньшей размерности сложнее описать особенности речи. Более крупные модели ($H=256$, 8 блоков, $N=256$) за 30 эпох уступили базовой. Почти у всех запусков качество росло до последней эпохи, поэтому более крупным моделям могло потребоваться больше времени.

\subsubsection{SepFormer}

В базовой конфигурации использовали $N=128$ фильтров длиной $L=16$ и чанки длиной $K=250$. Внутри чанков и между ними признаки обрабатывали 2 блока по 4 слоя трансформера с 8 головами в механизме attention и feed-forward (FF) размерности 512 (3.26\,M параметров). Обучали с AdamW и OneCycleLR (LR до $5\cdot10^{-4}$). Каждый запуск видел $192\,000$ примеров. В конфигурации из статьи~\cite{subakan2021sepformer} использовали $N=256$, 2 блока по 8 слоёв и FF размерности 1024.

\begin{experimenttable}
    \caption{Масштабирование SepFormer, 20 эпох}\label{tab:sepformer_scaling}
    \setlength{\tabcolsep}{1.8pt}
    \begin{tabularx}{\linewidth}{@{}>{\raggedright\arraybackslash}Xrrrrr@{}}
        \toprule
Изменение & Парам. & SI-SNRi & SDRi & PESQ & STOI \\
        \midrule
Базовая                & 3.26  & 4.76 & 5.22 & 1.33 & 0.759 \\
4 блока                & 6.43  & 5.24 & 5.70 & 1.36 & 0.773 \\
8 слоёв                & 6.43  & 5.42 & 5.87 & 1.37 & 0.779 \\
$N=192$, FF 768        & 7.31  & 5.53 & 5.98 & 1.38 & 0.781 \\
Из статьи              & 25.61 & \bestmetric{6.26} & \bestmetric{6.69} & \bestmetric{1.43} & \bestmetric{0.800} \\
        \bottomrule
    \end{tabularx}
\end{experimenttable}

При увеличении размера модели SI-SNRi рос и в конфигурации из статьи превысил базовый результат на $1.50$. При одинаковом числе параметров увеличение глубины блока до 8 слоёв дало больший прирост, чем добавление блоков. Сокращение чанка до $K=100$ снизило SI-SNRi до $4.62$, что согласуется с пользой более длинного контекста для SepFormer. Его отставание от остальных моделей более чем на 3~дБ может быть связано с коротким обучением и объёмом данных. Выборка из 20\,000 смесей длительностью 2~с содержит около 11 часов звука.

\subsubsection{Сравнение моделей}

В табл.~\ref{tab:audio_best} приведены результаты долгого обучения Conv-TasNet с $H=512$, DPRNN с $L=8$ и батчем 16 и базового SepFormer (54 эпохи для SepFormer). SI-SNRi$_{30}$ --- лучший результат за первые 30 эпох по логам обучения. Среди этих моделей DPRNN показала лучший SI-SNRi и за первые 30 эпох ($10.17$), и после 100 эпох ($12.48$). При 3.60\,M параметров она также достигла PESQ $2.17$ и STOI $0.907$.

\begin{experimenttable}
    \caption{Аудиомодели после долгого обучения}\label{tab:audio_best}
    \setlength{\tabcolsep}{1.8pt}
    \begin{tabularx}{\linewidth}{@{}>{\raggedright\arraybackslash}Xrrrrr@{}}
        \toprule
Модель & Парам. & SI-SNRi$_{30}$ & SI-SNRi & PESQ & STOI \\
        \midrule
Conv-TasNet     & 6.56 & 9.93 & 10.75 & 1.92 & 0.881 \\
DPRNN, $L=8$, батч 16 & 3.60 & \bestmetric{10.17} & \bestmetric{12.48} & \bestmetric{2.17} & \bestmetric{0.907} \\
SepFormer       & 3.26 & 6.16 & 8.01  & 1.59 & 0.838 \\
        \bottomrule
    \end{tabularx}
\end{experimenttable}


\subsection{Сравнение аудиомоделей и видеоэнкодеров}

В табл.~\ref{tab:main_comparison} Audio обозначает обучение только по звуку. Во всех 4 вариантах с видео использовали сложение признаков. Для DPRNN Audio взяты результаты базовой конфигурации из табл.~\ref{tab:dprnn_ablation}, где SI-SNR отдельно не приведён.

\begin{experimenttable}
    \caption{Разделение речи при обучении с нуля, 30 эпох, батч 32}\label{tab:main_comparison}
    \setlength{\tabcolsep}{1.8pt}
    \begin{tabularx}{\linewidth}{@{}>{\raggedright\arraybackslash}Xrrrrr@{}}
        \toprule
Модель & SI-SNR & SI-SNRi & SDRi & PESQ & STOI \\
        \midrule
Conv-TasNet Audio & 9.00 & 9.00 & 9.44 & 1.69 & 0.85 \\
DPRNN Audio & 9.83 & 9.83 & 10.28 & 1.78 & 0.86 \\
Conv-TasNet + Conv3D & 10.56 & 10.56 & 10.98 & 1.92 & 0.88 \\
Conv-TasNet + ViT & 10.32 & 10.32 & 10.74 & 1.88 & 0.88 \\
DPRNN + Conv3D & \bestmetric{11.37} & \bestmetric{11.37} & \bestmetric{11.78} & \bestmetric{1.98} & \bestmetric{0.89} \\
DPRNN + ViT & 11.31 & 11.31 & 11.72 & 1.97 & 0.89 \\
        \bottomrule
    \end{tabularx}
\end{experimenttable}

По SI-SNRi базовая DPRNN опередила Conv-TasNet на $0.83$. При добавлении Conv3D прирост SI-SNRi составил $1.56$ для Conv-TasNet и $1.54$ для DPRNN. Среди моделей в табл.~\ref{tab:main_comparison} лучшие значения всех 5 метрик показала DPRNN + Conv3D.

\subsection{Способ объединения звуковых и визуальных признаков}

В DPRNN + ViT сравнили сложение признаков (Add) и конкатенацию с MLP $192 \rightarrow 128 \rightarrow 128$ (Concat-MLP). Условия обучения совпадали с основным сравнением, функция потерь --- PIT SI-SNR.

\begin{experimenttable}
    \caption{Объединение признаков в DPRNN + ViT}\label{tab:fusion_comparison}
    \setlength{\tabcolsep}{1.8pt}
    \begin{tabularx}{\linewidth}{@{}>{\raggedright\arraybackslash}Xrrrrr@{}}
        \toprule
Метод & SI-SNR & SI-SNRi & SDRi & PESQ & STOI \\
        \midrule
Add & \bestmetric{11.31} & \bestmetric{11.31} & \bestmetric{11.72} & \bestmetric{1.97} & \bestmetric{0.89} \\
Concat-MLP & 9.10 & 9.10 & 9.51 & 1.73 & 0.86 \\
        \bottomrule
    \end{tabularx}
\end{experimenttable}

Сложение опередило Concat-MLP по всем 5 метрикам. Разница в SI-SNRi составила $2.21$, в PESQ~--- $0.25$.

\subsection{Стратегия обучения аудиовизуальной модели}

В режиме \textit{Scratch} аудио и видеочасть обучали вместе с нуля в течение 20 эпох. Для остальных режимов использовали один и тот же чекпоинт Conv-TasNet после 20 эпох обучения только на звуке. В \textit{Joint} сразу дообучали всю модель ещё 20 эпох. В \textit{Frozen} аудиочасть фиксировали на 5 эпох и обучали только видеоэнкодер и слой проекции. Это позволяло оценить, как меняется качество разделения за счёт обучения одной визуальной ветви. В \textit{Staged} после этих 5 эпох размораживали аудиочасть и дообучали всю модель ещё 15 эпох.

Во всех режимах размер батча составлял 16. Начальную скорость обучения задавали равной $10^{-3}$ для Scratch и Frozen и $3\cdot10^{-4}$ для Joint и этапа после разморозки.

\begin{experimenttable}
    \caption{Стратегии обучения Conv-TasNet + Conv3D}\label{tab:training_strategy}
    \setlength{\tabcolsep}{1.8pt}
    \begin{tabularx}{\linewidth}{@{}>{\raggedright\arraybackslash}Xrrrrr@{}}
        \toprule
Режим & SI-SNR & SI-SNRi & SDRi & PESQ & STOI \\
        \midrule
Scratch (с нуля) & \bestmetric{10.18} & \bestmetric{10.18} & \bestmetric{10.60} & \bestmetric{1.86} & \bestmetric{0.88} \\
Joint (совместное дообучение) & 9.62 & 9.62 & 10.06 & 1.78 & 0.86 \\
Frozen (звук заморожен) & 8.54 & 8.54 & 8.98 & 1.62 & 0.84 \\
Staged (заморозка и разморозка) & 9.55 & 9.55 & 9.98 & 1.76 & 0.86 \\
        \bottomrule
    \end{tabularx}
\end{experimenttable}

После разморозки SI-SNRi вырос с $8.54$ до $9.55$. Среди проверенных режимов Scratch показал лучшие значения всех 5 метрик. Сравнение с Joint учитывает также различие в начальном LR.

\subsection{Вспомогательная функция потерь}

К PIT SI-SNR для DPRNN + ViT (Add) добавляли cross энтропию с весом $\lambda$ для сопоставления видео с чистой речью. Бюджет обучения соответствовал основному сравнению.

\begin{experimenttable}
    \caption{DPRNN + ViT с разным весом CE, 30 эпох}\label{tab:auxiliary_loss}
    \setlength{\tabcolsep}{1.8pt}
    \begin{tabularx}{\linewidth}{@{}>{\raggedright\arraybackslash}Xrrrrr@{}}
        \toprule
Вес $\lambda$ & SI-SNR & SI-SNRi & SDRi & PESQ & STOI \\
        \midrule
Без CE & 11.31 & 11.31 & 11.72 & 1.97 & 0.89 \\
0.1 & \bestmetric{11.44} & \bestmetric{11.44} & \bestmetric{11.86} & \bestmetric{1.99} & \bestmetric{0.89} \\
0.5 & 11.15 & 11.15 & 11.55 & 1.97 & 0.89 \\
0.75 & 11.18 & 11.19 & 11.61 & 1.93 & 0.89 \\
        \bottomrule
    \end{tabularx}
\end{experimenttable}

При $\lambda=0.1$ все 5 метрик превысили результат DPRNN + ViT без CE. Прирост SI-SNRi составил $0.14$. При $\lambda=0.5$ и $0.75$ SI-SNRi оказался ниже базового.

\subsection{Длительное обучение итоговой модели}

За базовую модель взяли DPRNN + ViT со сложением признаков и вспомогательной функцией потерь CE при $\lambda=0.1$, обученную 30 эпох с батчем 32 и длиной ядра аудиоэнкодера 16 отсчётов. Для итогового запуска сохранили способ объединения признаков и функцию потерь, использовали ядро длиной 8 отсчётов, батч 16 и нормировку. Обучение длилось 100 эпох. В табл.~\ref{tab:long_training} приведены результаты на валидационной выборке.

\begin{experimenttable}
    \caption{Базовая и итоговая конфигурации DPRNN + ViT}
    \label{tab:long_training}
    \setlength{\tabcolsep}{1.8pt}
    \begin{tabularx}{\linewidth}{@{}>{\raggedright\arraybackslash}Xrrrrr@{}}
        \toprule
        Запуск & SI-SNR & SI-SNRi & SDRi & PESQ & STOI \\
        \midrule
        Базовый, 30 эпох & 11.44 & 11.44 & 11.86 & 1.99 & 0.89 \\
        Итоговый, 100 эпох & \bestmetric{12.69} & \bestmetric{12.69} & \bestmetric{13.09} & \bestmetric{2.21} & \bestmetric{0.91} \\
        \bottomrule
    \end{tabularx}
\end{experimenttable}

Лучший чекпоинт итогового запуска получен на 63-й эпохе. Он превзошёл базовую модель по всем 5 метрикам, а прирост SI-SNRi составил $1.25$.

\section{Заключение}

Мы разработали систему, которая восстанавливает 2 отдельные речевые дорожки из одноканальной записи одновременно говорящих людей. Сначала мы сравнили Conv-TasNet, DPRNN и SepFormer при разделении только по звуку. Затем расширили Conv-TasNet и DPRNN, чтобы при разделении смеси они могли учитывать движения губ обоих дикторов. При одинаковом бюджете в 30 эпох добавление свёрточного видеоэнкодера улучшило SI-SNRi у обеих моделей.

Лучшей среди аудиомоделей стала DPRNN с ядром отсчётов длиной 8 и SI-SNRi $12.48$ на валидационной выборке. Для итоговой аудиовизуальной системы мы выбрали DPRNN с энкодером ViT и сложением звуковых и визуальных признаков. Помимо разделения смеси модель училась сопоставлять движения губ каждого диктора с его голосом. Лучший чекпоинт этой системы достиг SI-SNRi $12.69$, PESQ $2.21$ и STOI $0.91$ на той же выборке, опередив лучшую аудиомодель по SI-SNRi на $0.21$.

\section*{Вклад авторов}
\noindent\textbf{Никита Романенко} --- реализовал аудиомодели Conv-TasNet, DPRNN и SepFormer, провёл эксперименты с ними (раздел~\ref{sec:audio_experiments}), а также написал инференс моделей, подсчёт метрик, замер эффективности и демо-ноутбук.\par

\noindent\textbf{Андрей Теплов} --- реализовал видеоэнкодеры Conv3D и ViT и их интеграцию в Conv-TasNet и DPRNN, добавил поддержку обучения аудио и аудиовизуальных моделей, провёл эксперименты с видео.\par

\bibliographystyle{IEEEtran}
\bibliography{references}

\end{document}
